% GENERATED BY src/python/lmi/tables/ext_asyl_employment.py
% -- DO NOT EDIT
%==============================================================================%
% ext_asyl_employment.tex
%
% Purpose : Erweiterung, Tabelle E1 -- Kursangebot x Asylstatus auf die
%           Erwerbstaetigkeit. Drei Tabellen: Koeffizienten, marginale Effekte
%           je Statusgruppe, sowie der Vergleich mit Tabelle S6 Modell 1.1.9
%           der Publikation, die dieser Job mitrepliziert.
% Source  : SOEPremote job 244835, src/stata/remote/ext_asyl_employment.do
%           Rohlisting: results/transcripts/
%                       src_stata_remote_ext_asyl_employment.do.eml
%           Transkript: results/comparisons/ext_asyl_employment.md
% Needs   : booktabs, tabularx (beide bereits in bachelorarbeit.sty geladen)
% Usage   : % GENERATED BY src/python/lmi/tables/ext_asyl_employment.py
% -- DO NOT EDIT
%==============================================================================%
% ext_asyl_employment.tex
%
% Purpose : Erweiterung, Tabelle E1 -- Kursangebot x Asylstatus auf die
%           Erwerbstaetigkeit. Drei Tabellen: Koeffizienten, marginale Effekte
%           je Statusgruppe, sowie der Vergleich mit Tabelle S6 Modell 1.1.9
%           der Publikation, die dieser Job mitrepliziert.
% Source  : SOEPremote job 244835, src/stata/remote/ext_asyl_employment.do
%           Rohlisting: results/transcripts/
%                       src_stata_remote_ext_asyl_employment.do.eml
%           Transkript: results/comparisons/ext_asyl_employment.md
% Needs   : booktabs, tabularx (beide bereits in bachelorarbeit.sty geladen)
% Usage   : % GENERATED BY src/python/lmi/tables/ext_asyl_employment.py
% -- DO NOT EDIT
%==============================================================================%
% ext_asyl_employment.tex
%
% Purpose : Erweiterung, Tabelle E1 -- Kursangebot x Asylstatus auf die
%           Erwerbstaetigkeit. Drei Tabellen: Koeffizienten, marginale Effekte
%           je Statusgruppe, sowie der Vergleich mit Tabelle S6 Modell 1.1.9
%           der Publikation, die dieser Job mitrepliziert.
% Source  : SOEPremote job 244835, src/stata/remote/ext_asyl_employment.do
%           Rohlisting: results/transcripts/
%                       src_stata_remote_ext_asyl_employment.do.eml
%           Transkript: results/comparisons/ext_asyl_employment.md
% Needs   : booktabs, tabularx (beide bereits in bachelorarbeit.sty geladen)
% Usage   : % GENERATED BY src/python/lmi/tables/ext_asyl_employment.py
% -- DO NOT EDIT
%==============================================================================%
% ext_asyl_employment.tex
%
% Purpose : Erweiterung, Tabelle E1 -- Kursangebot x Asylstatus auf die
%           Erwerbstaetigkeit. Drei Tabellen: Koeffizienten, marginale Effekte
%           je Statusgruppe, sowie der Vergleich mit Tabelle S6 Modell 1.1.9
%           der Publikation, die dieser Job mitrepliziert.
% Source  : SOEPremote job 244835, src/stata/remote/ext_asyl_employment.do
%           Rohlisting: results/transcripts/
%                       src_stata_remote_ext_asyl_employment.do.eml
%           Transkript: results/comparisons/ext_asyl_employment.md
% Needs   : booktabs, tabularx (beide bereits in bachelorarbeit.sty geladen)
% Usage   : \input{tables/ext_asyl_employment}
% Author  : Emre Suemer, 2026-09-01
%==============================================================================%

\makeatletter
\@ifundefined{NC@rewrite@Z}{%
  \newcolumntype{Z}{>{\raggedright\arraybackslash\hangindent=1.9em\hangafter=1}X}%
}{}
\makeatother

\begin{table}[htbp]
  \centering
  \caption{Kursangebot und Asylstatus: Wahrscheinlichkeit bezahlter Arbeit}
  \label{tab:ext-asyl-erwerb}
  \footnotesize
  \setlength{\tabcolsep}{4pt}%
  \begin{tabularx}{\textwidth}{@{}Zrr@{}}
    \toprule
     & Modell E1.1 & Modell E1.2 \\
     & Status im & Status bei \\
     & Erhebungsjahr & Erstbeobachtung \\
     & p.p. (SE) & p.p. (SE) \\
    \midrule
    \textbf{Kreisangebot an Integrationskursen, std.} & 3{,}97$^{*}$ (1{,}55) & 4{,}25$^{**}$ (1{,}31) \\
    \quad $\times$ Anerkannt & $-$1{,}28 (1{,}90) & $-$2{,}87 (2{,}03) \\
    \quad $\times$ Abgelehnt & $-$2{,}56 (2{,}47) & $-$5{,}13$^{*}$ (2{,}55) \\
    Aufenthaltsdauer, in Monaten (/12) & 10{,}73$^{**}$ (1{,}60) & 11{,}37$^{**}$ (1{,}60) \\
    Bildungsjahre vor der Migration & 0{,}56$^{**}$ (0{,}19) & 0{,}57$^{**}$ (0{,}20) \\
    Erwerbserfahrung vor der Migration & 4{,}64$^{+}$ (2{,}38) & 4{,}56$^{+}$ (2{,}39) \\
    Weiblich & $-$15{,}76$^{**}$ (2{,}16) & $-$15{,}69$^{**}$ (2{,}18) \\
    Kinder unter 16 Jahren im Haushalt & $-$11{,}02$^{**}$ (2{,}15) & $-$11{,}27$^{**}$ (2{,}18) \\
    Einreise mit der Familie & 5{,}33$^{**}$ (1{,}97) & 5{,}38$^{**}$ (2{,}01) \\
    Unterst\"utzung durch Familie/Verwandte in DE & 4{,}01 (2{,}83) & 4{,}10 (2{,}90) \\
    Alter bei der Erstbefragung & $-$0{,}35$^{**}$ (0{,}09) & $-$0{,}34$^{**}$ (0{,}09) \\
    Traumatische Erfahrungen & $-$0{,}24 (2{,}31) & $-$0{,}11 (2{,}33) \\
    Status des Asylantrags (Ref.: keine Entscheidung) & & \\
    \quad -- Anerkannt & 5{,}27$^{*}$ (2{,}61) & 5{,}21$^{*}$ (2{,}63) \\
    \quad -- Abgelehnt & 8{,}20$^{**}$ (2{,}76) & 4{,}39 (2{,}95) \\
    Kreisanteil Ausl\"ander:innen, std. & $-$1{,}12 (2{,}79) & $-$0{,}80 (2{,}83) \\
    Kreisbev\"olkerungsdichte, std. & $-$0{,}95 (3{,}61) & $-$1{,}05 (3{,}62) \\
    Kreisarbeitslosenquote, std. & $-$4{,}43$^{+}$ (2{,}43) & $-$3{,}99 (2{,}48) \\
    Kreissorgen \"uber Zuwanderung, std. & 0{,}33 (1{,}19) & 0{,}30 (1{,}21) \\
    Kreisanteil CDU/CSU-W\"ahler:innen, std. & $-$2{,}01 (2{,}20) & $-$1{,}65 (2{,}21) \\
    Konstante & $-$2{,}82 (6{,}31) & $-$4{,}16 (6{,}36) \\
    \midrule
    \multicolumn{3}{@{}l}{\textit{Fixe Effekte}} \\
    \quad Erhebungsjahr, Erhebungsstichprobe & Ja & Ja \\
    \quad Bundesland, Herkunftsland (aggregiert) & Ja & Ja \\
    \midrule
    Gemeinsamer Test beider Interaktionen & F(2; 2726{,}6)\,=\,0{,}57 & F(2; 2726{,}9)\,=\,2{,}48 \\
    \quad p & 0{,}563 & 0{,}084 \\
    \midrule
    N Beobachtungen & 5.467 & 5.467 \\
    N Personen & 2.730 & 2.730 \\
    N Imputationen & 20 & 20 \\
    R\textsuperscript{2} (adj.) & 0{,}2137 & 0{,}2121 \\
    \bottomrule
  \end{tabularx}

  \vspace{0.75ex}
  \parbox{\textwidth}{\footnotesize
    \textit{Anmerkungen:} Eigene Berechnung, IAB-BAMF-SOEP-Befragung von
    Gefl\"uchteten 2016--2019, SOEP-Core v36 (Remote Edition), gewichtet.  $^{+}$~p\,$<$\,0{,}1, $^{*}$~p\,$<$\,0{,}05, $^{**}$~p\,$<$\,0{,}01
    (zweiseitig).}
    
	%   Spezifikation wie Modell~1.1 in Tabelle~\ref{tab:erwerbstaetigkeit},
	%   erg\"anzt um die Interaktion des Kursangebots mit dem Asylstatus.
	%   Lineare Wahrscheinlichkeitsmodelle, durchschnittliche marginale Effekte in
	%   Prozentpunkten; robuste, auf Personenebene geclusterte Standardfehler in
	%   Klammern. Referenzkategorie des Asylstatus ist \glqq keine
	%   Entscheidung\grqq{}; die Kategorien sind gegen\"uber der Publikation
	%   korrigiert (vgl.\ Tabelle~\ref{tab:erwerbstaetigkeit}).
	%   Modell~E1.1 misst den Status im jeweiligen Erhebungsjahr, Modell~E1.2 bei
	%   der Erstbeobachtung der Person -- Letzteres ist konzeptionell sauberer,
	%   weil Asylentscheidungen erst \emph{nach} der Kreiszuweisung ergehen
	%   (Tabelle~\ref{tab:ext-asyl-uebergang}). In Spalte~E1.2 beziehen sich
	%   sowohl die Interaktions- als auch die Haupteffektzeilen des Asylstatus
	%   entsprechend auf den Status bei Erstbeobachtung.
    
\end{table}

\begin{table}[htbp]
  \centering
  \caption{Marginaler Effekt des Kursangebots auf die Erwerbst\"atigkeit,
           nach Asylstatus}
  \label{tab:ext-asyl-erwerb-ame}
  \footnotesize
  \setlength{\tabcolsep}{4pt}%
  \begin{tabularx}{\textwidth}{@{}Zrrrr@{}}
    \toprule
     & \multicolumn{2}{c}{Modell E1.1: Erhebungsjahr}
     & \multicolumn{2}{c}{Modell E1.2: Erstbeobachtung} \\
    \cmidrule(lr){2-3}\cmidrule(lr){4-5}
    Status des Asylantrags & p.p. (SE) & p & p.p. (SE) & p \\
    \midrule
    Keine Entscheidung & 4{,}0$^{*}$ (1{,}5) & 0{,}010 & 4{,}3$^{**}$ (1{,}3) & 0{,}001 \\
    Anerkannt & 2{,}7$^{+}$ (1{,}5) & 0{,}078 & 1{,}4 (1{,}9) & 0{,}467 \\
    Abgelehnt & 1{,}4 (2{,}2) & 0{,}519 & $-$0{,}9 (2{,}4) & 0{,}720 \\
    \midrule
    Gemeinsamer Test & \multicolumn{2}{c}{F(2; 2726{,}6)\,=\,0{,}57, p\,=\,0{,}563}
     & \multicolumn{2}{c}{F(2; 2726{,}9)\,=\,2{,}48, p\,=\,0{,}084} \\
    \bottomrule
  \end{tabularx}

  \vspace{0.75ex}
  \parbox{\textwidth}{\footnotesize
    \textit{Anmerkungen:} Durchschnittliche marginale Effekte einer Erh\"ohung
    des Kreisangebots um eine Standardabweichung auf die Wahrscheinlichkeit
    bezahlter Arbeit, in Prozentpunkten, aus den Modellen der
    Tabelle~\ref{tab:ext-asyl-erwerb} (\texttt{mimrgns}, Poolung nach Rubins
    Regeln). Kontrolliert ist die Identit\"at
    AME(Gruppe)\,=\,$b$[Angebot]\,$+$\,$b$[Angebot\,$\times$\,Gruppe];
    gr\"o\ss{}te Abweichung zwischen beiden Rechenwegen 0{,}050~p.p.
    \\[0.4ex]
    \textbf{Inferenzregel, vorab festgelegt} (Designdokument, Abschnitt~3e):
    Prim\"ar ist der gemeinsame Test beider Interaktionsterme, einzelne
    Kontraste sind deskriptiv. \emph{Keine} der beiden Spezifikationen
    erreicht p\,$<$\,0{,}05. Der einzelne signifikante Kontrast in
    Modell~E1.2 darf daher nicht als Befund berichtet werden.
    \\[0.4ex]
    \textbf{Statistische Power, ebenfalls vorab festgehalten:} Bei
    Gruppenanteilen von 39/38/23~Prozent liegen die Standardfehler der
    Subgruppen etwa beim 1{,}6- bis 2{,}1-Fachen des gepoolten Werts von
    1{,}19. Ein Unterschied unter rund 4~Prozentpunkten ist auf dieser
    Ergebnisstufe nicht nachweisbar. Die Aussage lautet deshalb nicht
    \glqq es gibt keinen Effekt\grqq{}, sondern: ein Effekt dieser
    Gr\"o\ss{}enordnung h\"atte hier nicht entdeckt werden k\"onnen.
    Die Erweiterung wird auf der Mechanismenstufe entschieden
    (Tabelle~\ref{tab:ext-asyl-mechanismen}).}
\end{table}

\begin{table}[htbp]
  \centering
  \caption{Modell~E1.1 im Vergleich zu Tabelle~S6, Modell~1.1.9 der
           Publikation}
  \label{tab:ext-asyl-s6}
  \footnotesize
  \setlength{\tabcolsep}{4pt}%
  \begin{tabularx}{\textwidth}{@{}Zrrr@{}}
    \toprule
     & Publikation & Replikation & Abweichung \\
    \midrule
    Kreisangebot, std. (d.\,h.\ f\"ur \glqq keine Entscheidung\grqq{}) & 3{,}98$^{*}$ (1{,}55) & 3{,}97$^{*}$ (1{,}55) & 0{,}01 \\
    \quad $\times$ Anerkannt & $-$1{,}30 (1{,}90) & $-$1{,}28 (1{,}90) & 0{,}02 \\
    \quad $\times$ Abgelehnt & $-$2{,}58 (2{,}48) & $-$2{,}56 (2{,}47) & 0{,}02 \\
    \midrule
    N Beobachtungen & 5.473 & 5.467 & 6 \\
    N Personen & 2.732 & 2.730 & 2 \\
    \bottomrule
  \end{tabularx}

  \vspace{0.75ex}
  \parbox{\textwidth}{\footnotesize
    \textit{Anmerkungen:} Publizierte Werte aus dem Zusatzmaterial zu
    \textcite{Kanas:2023}, S.~12, Tabelle~S6, Modell~1.1.9; eigene Werte aus
    SOEPremote-Job~244835 (Modell~E1.1). Die Autorinnen sch\"atzen dort exakt
    diese Spezifikation, weshalb Modell~E1.1 zugleich eine Replikation von
    S6 ist. Alle drei Koeffizienten stimmen bis auf 0{,}02~p.p.
    \"uberein, die Standardfehler bis auf 0{,}01, und die Signifikanzniveaus
    vollst\"andig -- dieselbe Toleranz wie in den f\"unf Haupttabellen; die
    Differenz ist Monte-Carlo-Variation der multiplen Imputation.
    \\[0.4ex]
    \textbf{Vierter Fehler in der Publikation:} Tabelle~S6 weist
    5.473 Beobachtungen und 2.732 Personen aus, w\"ahrend die Tabellen~3
    bis~5 auf derselben Stichprobe und mit denselben Kontrollvariablen
    5.467 und 2.730 ausweisen. Die Nachsch\"atzung ergibt 5.467/2.730,
    passend zum Stichprobentrichter
    (Tabelle~\ref{tab:stichprobenselektion}). Die sechs zus\"atzlichen
    Personenjahre und zwei zus\"atzlichen Personen in S6 existieren nicht.
    Die Sch\"atzwerte sind davon unber\"uhrt.}
\end{table}

% Author  : Emre Suemer, 2026-09-01
%==============================================================================%

\makeatletter
\@ifundefined{NC@rewrite@Z}{%
  \newcolumntype{Z}{>{\raggedright\arraybackslash\hangindent=1.9em\hangafter=1}X}%
}{}
\makeatother

\begin{table}[htbp]
  \centering
  \caption{Kursangebot und Asylstatus: Wahrscheinlichkeit bezahlter Arbeit}
  \label{tab:ext-asyl-erwerb}
  \footnotesize
  \setlength{\tabcolsep}{4pt}%
  \begin{tabularx}{\textwidth}{@{}Zrr@{}}
    \toprule
     & Modell E1.1 & Modell E1.2 \\
     & Status im & Status bei \\
     & Erhebungsjahr & Erstbeobachtung \\
     & p.p. (SE) & p.p. (SE) \\
    \midrule
    \textbf{Kreisangebot an Integrationskursen, std.} & 3{,}97$^{*}$ (1{,}55) & 4{,}25$^{**}$ (1{,}31) \\
    \quad $\times$ Anerkannt & $-$1{,}28 (1{,}90) & $-$2{,}87 (2{,}03) \\
    \quad $\times$ Abgelehnt & $-$2{,}56 (2{,}47) & $-$5{,}13$^{*}$ (2{,}55) \\
    Aufenthaltsdauer, in Monaten (/12) & 10{,}73$^{**}$ (1{,}60) & 11{,}37$^{**}$ (1{,}60) \\
    Bildungsjahre vor der Migration & 0{,}56$^{**}$ (0{,}19) & 0{,}57$^{**}$ (0{,}20) \\
    Erwerbserfahrung vor der Migration & 4{,}64$^{+}$ (2{,}38) & 4{,}56$^{+}$ (2{,}39) \\
    Weiblich & $-$15{,}76$^{**}$ (2{,}16) & $-$15{,}69$^{**}$ (2{,}18) \\
    Kinder unter 16 Jahren im Haushalt & $-$11{,}02$^{**}$ (2{,}15) & $-$11{,}27$^{**}$ (2{,}18) \\
    Einreise mit der Familie & 5{,}33$^{**}$ (1{,}97) & 5{,}38$^{**}$ (2{,}01) \\
    Unterst\"utzung durch Familie/Verwandte in DE & 4{,}01 (2{,}83) & 4{,}10 (2{,}90) \\
    Alter bei der Erstbefragung & $-$0{,}35$^{**}$ (0{,}09) & $-$0{,}34$^{**}$ (0{,}09) \\
    Traumatische Erfahrungen & $-$0{,}24 (2{,}31) & $-$0{,}11 (2{,}33) \\
    Status des Asylantrags (Ref.: keine Entscheidung) & & \\
    \quad -- Anerkannt & 5{,}27$^{*}$ (2{,}61) & 5{,}21$^{*}$ (2{,}63) \\
    \quad -- Abgelehnt & 8{,}20$^{**}$ (2{,}76) & 4{,}39 (2{,}95) \\
    Kreisanteil Ausl\"ander:innen, std. & $-$1{,}12 (2{,}79) & $-$0{,}80 (2{,}83) \\
    Kreisbev\"olkerungsdichte, std. & $-$0{,}95 (3{,}61) & $-$1{,}05 (3{,}62) \\
    Kreisarbeitslosenquote, std. & $-$4{,}43$^{+}$ (2{,}43) & $-$3{,}99 (2{,}48) \\
    Kreissorgen \"uber Zuwanderung, std. & 0{,}33 (1{,}19) & 0{,}30 (1{,}21) \\
    Kreisanteil CDU/CSU-W\"ahler:innen, std. & $-$2{,}01 (2{,}20) & $-$1{,}65 (2{,}21) \\
    Konstante & $-$2{,}82 (6{,}31) & $-$4{,}16 (6{,}36) \\
    \midrule
    \multicolumn{3}{@{}l}{\textit{Fixe Effekte}} \\
    \quad Erhebungsjahr, Erhebungsstichprobe & Ja & Ja \\
    \quad Bundesland, Herkunftsland (aggregiert) & Ja & Ja \\
    \midrule
    Gemeinsamer Test beider Interaktionen & F(2; 2726{,}6)\,=\,0{,}57 & F(2; 2726{,}9)\,=\,2{,}48 \\
    \quad p & 0{,}563 & 0{,}084 \\
    \midrule
    N Beobachtungen & 5.467 & 5.467 \\
    N Personen & 2.730 & 2.730 \\
    N Imputationen & 20 & 20 \\
    R\textsuperscript{2} (adj.) & 0{,}2137 & 0{,}2121 \\
    \bottomrule
  \end{tabularx}

  \vspace{0.75ex}
  \parbox{\textwidth}{\footnotesize
    \textit{Anmerkungen:} Eigene Berechnung, IAB-BAMF-SOEP-Befragung von
    Gefl\"uchteten 2016--2019, SOEP-Core v36 (Remote Edition), gewichtet.  $^{+}$~p\,$<$\,0{,}1, $^{*}$~p\,$<$\,0{,}05, $^{**}$~p\,$<$\,0{,}01
    (zweiseitig).}
    
	%   Spezifikation wie Modell~1.1 in Tabelle~\ref{tab:erwerbstaetigkeit},
	%   erg\"anzt um die Interaktion des Kursangebots mit dem Asylstatus.
	%   Lineare Wahrscheinlichkeitsmodelle, durchschnittliche marginale Effekte in
	%   Prozentpunkten; robuste, auf Personenebene geclusterte Standardfehler in
	%   Klammern. Referenzkategorie des Asylstatus ist \glqq keine
	%   Entscheidung\grqq{}; die Kategorien sind gegen\"uber der Publikation
	%   korrigiert (vgl.\ Tabelle~\ref{tab:erwerbstaetigkeit}).
	%   Modell~E1.1 misst den Status im jeweiligen Erhebungsjahr, Modell~E1.2 bei
	%   der Erstbeobachtung der Person -- Letzteres ist konzeptionell sauberer,
	%   weil Asylentscheidungen erst \emph{nach} der Kreiszuweisung ergehen
	%   (Tabelle~\ref{tab:ext-asyl-uebergang}). In Spalte~E1.2 beziehen sich
	%   sowohl die Interaktions- als auch die Haupteffektzeilen des Asylstatus
	%   entsprechend auf den Status bei Erstbeobachtung.
    
\end{table}

\begin{table}[htbp]
  \centering
  \caption{Marginaler Effekt des Kursangebots auf die Erwerbst\"atigkeit,
           nach Asylstatus}
  \label{tab:ext-asyl-erwerb-ame}
  \footnotesize
  \setlength{\tabcolsep}{4pt}%
  \begin{tabularx}{\textwidth}{@{}Zrrrr@{}}
    \toprule
     & \multicolumn{2}{c}{Modell E1.1: Erhebungsjahr}
     & \multicolumn{2}{c}{Modell E1.2: Erstbeobachtung} \\
    \cmidrule(lr){2-3}\cmidrule(lr){4-5}
    Status des Asylantrags & p.p. (SE) & p & p.p. (SE) & p \\
    \midrule
    Keine Entscheidung & 4{,}0$^{*}$ (1{,}5) & 0{,}010 & 4{,}3$^{**}$ (1{,}3) & 0{,}001 \\
    Anerkannt & 2{,}7$^{+}$ (1{,}5) & 0{,}078 & 1{,}4 (1{,}9) & 0{,}467 \\
    Abgelehnt & 1{,}4 (2{,}2) & 0{,}519 & $-$0{,}9 (2{,}4) & 0{,}720 \\
    \midrule
    Gemeinsamer Test & \multicolumn{2}{c}{F(2; 2726{,}6)\,=\,0{,}57, p\,=\,0{,}563}
     & \multicolumn{2}{c}{F(2; 2726{,}9)\,=\,2{,}48, p\,=\,0{,}084} \\
    \bottomrule
  \end{tabularx}

  \vspace{0.75ex}
  \parbox{\textwidth}{\footnotesize
    \textit{Anmerkungen:} Durchschnittliche marginale Effekte einer Erh\"ohung
    des Kreisangebots um eine Standardabweichung auf die Wahrscheinlichkeit
    bezahlter Arbeit, in Prozentpunkten, aus den Modellen der
    Tabelle~\ref{tab:ext-asyl-erwerb} (\texttt{mimrgns}, Poolung nach Rubins
    Regeln). Kontrolliert ist die Identit\"at
    AME(Gruppe)\,=\,$b$[Angebot]\,$+$\,$b$[Angebot\,$\times$\,Gruppe];
    gr\"o\ss{}te Abweichung zwischen beiden Rechenwegen 0{,}050~p.p.
    \\[0.4ex]
    \textbf{Inferenzregel, vorab festgelegt} (Designdokument, Abschnitt~3e):
    Prim\"ar ist der gemeinsame Test beider Interaktionsterme, einzelne
    Kontraste sind deskriptiv. \emph{Keine} der beiden Spezifikationen
    erreicht p\,$<$\,0{,}05. Der einzelne signifikante Kontrast in
    Modell~E1.2 darf daher nicht als Befund berichtet werden.
    \\[0.4ex]
    \textbf{Statistische Power, ebenfalls vorab festgehalten:} Bei
    Gruppenanteilen von 39/38/23~Prozent liegen die Standardfehler der
    Subgruppen etwa beim 1{,}6- bis 2{,}1-Fachen des gepoolten Werts von
    1{,}19. Ein Unterschied unter rund 4~Prozentpunkten ist auf dieser
    Ergebnisstufe nicht nachweisbar. Die Aussage lautet deshalb nicht
    \glqq es gibt keinen Effekt\grqq{}, sondern: ein Effekt dieser
    Gr\"o\ss{}enordnung h\"atte hier nicht entdeckt werden k\"onnen.
    Die Erweiterung wird auf der Mechanismenstufe entschieden
    (Tabelle~\ref{tab:ext-asyl-mechanismen}).}
\end{table}

\begin{table}[htbp]
  \centering
  \caption{Modell~E1.1 im Vergleich zu Tabelle~S6, Modell~1.1.9 der
           Publikation}
  \label{tab:ext-asyl-s6}
  \footnotesize
  \setlength{\tabcolsep}{4pt}%
  \begin{tabularx}{\textwidth}{@{}Zrrr@{}}
    \toprule
     & Publikation & Replikation & Abweichung \\
    \midrule
    Kreisangebot, std. (d.\,h.\ f\"ur \glqq keine Entscheidung\grqq{}) & 3{,}98$^{*}$ (1{,}55) & 3{,}97$^{*}$ (1{,}55) & 0{,}01 \\
    \quad $\times$ Anerkannt & $-$1{,}30 (1{,}90) & $-$1{,}28 (1{,}90) & 0{,}02 \\
    \quad $\times$ Abgelehnt & $-$2{,}58 (2{,}48) & $-$2{,}56 (2{,}47) & 0{,}02 \\
    \midrule
    N Beobachtungen & 5.473 & 5.467 & 6 \\
    N Personen & 2.732 & 2.730 & 2 \\
    \bottomrule
  \end{tabularx}

  \vspace{0.75ex}
  \parbox{\textwidth}{\footnotesize
    \textit{Anmerkungen:} Publizierte Werte aus dem Zusatzmaterial zu
    \textcite{Kanas:2023}, S.~12, Tabelle~S6, Modell~1.1.9; eigene Werte aus
    SOEPremote-Job~244835 (Modell~E1.1). Die Autorinnen sch\"atzen dort exakt
    diese Spezifikation, weshalb Modell~E1.1 zugleich eine Replikation von
    S6 ist. Alle drei Koeffizienten stimmen bis auf 0{,}02~p.p.
    \"uberein, die Standardfehler bis auf 0{,}01, und die Signifikanzniveaus
    vollst\"andig -- dieselbe Toleranz wie in den f\"unf Haupttabellen; die
    Differenz ist Monte-Carlo-Variation der multiplen Imputation.
    \\[0.4ex]
    \textbf{Vierter Fehler in der Publikation:} Tabelle~S6 weist
    5.473 Beobachtungen und 2.732 Personen aus, w\"ahrend die Tabellen~3
    bis~5 auf derselben Stichprobe und mit denselben Kontrollvariablen
    5.467 und 2.730 ausweisen. Die Nachsch\"atzung ergibt 5.467/2.730,
    passend zum Stichprobentrichter
    (Tabelle~\ref{tab:stichprobenselektion}). Die sechs zus\"atzlichen
    Personenjahre und zwei zus\"atzlichen Personen in S6 existieren nicht.
    Die Sch\"atzwerte sind davon unber\"uhrt.}
\end{table}

% Author  : Emre Suemer, 2026-09-01
%==============================================================================%

\makeatletter
\@ifundefined{NC@rewrite@Z}{%
  \newcolumntype{Z}{>{\raggedright\arraybackslash\hangindent=1.9em\hangafter=1}X}%
}{}
\makeatother

\begin{table}[htbp]
  \centering
  \caption{Kursangebot und Asylstatus: Wahrscheinlichkeit bezahlter Arbeit}
  \label{tab:ext-asyl-erwerb}
  \footnotesize
  \setlength{\tabcolsep}{4pt}%
  \begin{tabularx}{\textwidth}{@{}Zrr@{}}
    \toprule
     & Modell E1.1 & Modell E1.2 \\
     & Status im & Status bei \\
     & Erhebungsjahr & Erstbeobachtung \\
     & p.p. (SE) & p.p. (SE) \\
    \midrule
    \textbf{Kreisangebot an Integrationskursen, std.} & 3{,}97$^{*}$ (1{,}55) & 4{,}25$^{**}$ (1{,}31) \\
    \quad $\times$ Anerkannt & $-$1{,}28 (1{,}90) & $-$2{,}87 (2{,}03) \\
    \quad $\times$ Abgelehnt & $-$2{,}56 (2{,}47) & $-$5{,}13$^{*}$ (2{,}55) \\
    Aufenthaltsdauer, in Monaten (/12) & 10{,}73$^{**}$ (1{,}60) & 11{,}37$^{**}$ (1{,}60) \\
    Bildungsjahre vor der Migration & 0{,}56$^{**}$ (0{,}19) & 0{,}57$^{**}$ (0{,}20) \\
    Erwerbserfahrung vor der Migration & 4{,}64$^{+}$ (2{,}38) & 4{,}56$^{+}$ (2{,}39) \\
    Weiblich & $-$15{,}76$^{**}$ (2{,}16) & $-$15{,}69$^{**}$ (2{,}18) \\
    Kinder unter 16 Jahren im Haushalt & $-$11{,}02$^{**}$ (2{,}15) & $-$11{,}27$^{**}$ (2{,}18) \\
    Einreise mit der Familie & 5{,}33$^{**}$ (1{,}97) & 5{,}38$^{**}$ (2{,}01) \\
    Unterst\"utzung durch Familie/Verwandte in DE & 4{,}01 (2{,}83) & 4{,}10 (2{,}90) \\
    Alter bei der Erstbefragung & $-$0{,}35$^{**}$ (0{,}09) & $-$0{,}34$^{**}$ (0{,}09) \\
    Traumatische Erfahrungen & $-$0{,}24 (2{,}31) & $-$0{,}11 (2{,}33) \\
    Status des Asylantrags (Ref.: keine Entscheidung) & & \\
    \quad -- Anerkannt & 5{,}27$^{*}$ (2{,}61) & 5{,}21$^{*}$ (2{,}63) \\
    \quad -- Abgelehnt & 8{,}20$^{**}$ (2{,}76) & 4{,}39 (2{,}95) \\
    Kreisanteil Ausl\"ander:innen, std. & $-$1{,}12 (2{,}79) & $-$0{,}80 (2{,}83) \\
    Kreisbev\"olkerungsdichte, std. & $-$0{,}95 (3{,}61) & $-$1{,}05 (3{,}62) \\
    Kreisarbeitslosenquote, std. & $-$4{,}43$^{+}$ (2{,}43) & $-$3{,}99 (2{,}48) \\
    Kreissorgen \"uber Zuwanderung, std. & 0{,}33 (1{,}19) & 0{,}30 (1{,}21) \\
    Kreisanteil CDU/CSU-W\"ahler:innen, std. & $-$2{,}01 (2{,}20) & $-$1{,}65 (2{,}21) \\
    Konstante & $-$2{,}82 (6{,}31) & $-$4{,}16 (6{,}36) \\
    \midrule
    \multicolumn{3}{@{}l}{\textit{Fixe Effekte}} \\
    \quad Erhebungsjahr, Erhebungsstichprobe & Ja & Ja \\
    \quad Bundesland, Herkunftsland (aggregiert) & Ja & Ja \\
    \midrule
    Gemeinsamer Test beider Interaktionen & F(2; 2726{,}6)\,=\,0{,}57 & F(2; 2726{,}9)\,=\,2{,}48 \\
    \quad p & 0{,}563 & 0{,}084 \\
    \midrule
    N Beobachtungen & 5.467 & 5.467 \\
    N Personen & 2.730 & 2.730 \\
    N Imputationen & 20 & 20 \\
    R\textsuperscript{2} (adj.) & 0{,}2137 & 0{,}2121 \\
    \bottomrule
  \end{tabularx}

  \vspace{0.75ex}
  \parbox{\textwidth}{\footnotesize
    \textit{Anmerkungen:} Eigene Berechnung, IAB-BAMF-SOEP-Befragung von
    Gefl\"uchteten 2016--2019, SOEP-Core v36 (Remote Edition), gewichtet.  $^{+}$~p\,$<$\,0{,}1, $^{*}$~p\,$<$\,0{,}05, $^{**}$~p\,$<$\,0{,}01
    (zweiseitig).}
    
	%   Spezifikation wie Modell~1.1 in Tabelle~\ref{tab:erwerbstaetigkeit},
	%   erg\"anzt um die Interaktion des Kursangebots mit dem Asylstatus.
	%   Lineare Wahrscheinlichkeitsmodelle, durchschnittliche marginale Effekte in
	%   Prozentpunkten; robuste, auf Personenebene geclusterte Standardfehler in
	%   Klammern. Referenzkategorie des Asylstatus ist \glqq keine
	%   Entscheidung\grqq{}; die Kategorien sind gegen\"uber der Publikation
	%   korrigiert (vgl.\ Tabelle~\ref{tab:erwerbstaetigkeit}).
	%   Modell~E1.1 misst den Status im jeweiligen Erhebungsjahr, Modell~E1.2 bei
	%   der Erstbeobachtung der Person -- Letzteres ist konzeptionell sauberer,
	%   weil Asylentscheidungen erst \emph{nach} der Kreiszuweisung ergehen
	%   (Tabelle~\ref{tab:ext-asyl-uebergang}). In Spalte~E1.2 beziehen sich
	%   sowohl die Interaktions- als auch die Haupteffektzeilen des Asylstatus
	%   entsprechend auf den Status bei Erstbeobachtung.
    
\end{table}

\begin{table}[htbp]
  \centering
  \caption{Marginaler Effekt des Kursangebots auf die Erwerbst\"atigkeit,
           nach Asylstatus}
  \label{tab:ext-asyl-erwerb-ame}
  \footnotesize
  \setlength{\tabcolsep}{4pt}%
  \begin{tabularx}{\textwidth}{@{}Zrrrr@{}}
    \toprule
     & \multicolumn{2}{c}{Modell E1.1: Erhebungsjahr}
     & \multicolumn{2}{c}{Modell E1.2: Erstbeobachtung} \\
    \cmidrule(lr){2-3}\cmidrule(lr){4-5}
    Status des Asylantrags & p.p. (SE) & p & p.p. (SE) & p \\
    \midrule
    Keine Entscheidung & 4{,}0$^{*}$ (1{,}5) & 0{,}010 & 4{,}3$^{**}$ (1{,}3) & 0{,}001 \\
    Anerkannt & 2{,}7$^{+}$ (1{,}5) & 0{,}078 & 1{,}4 (1{,}9) & 0{,}467 \\
    Abgelehnt & 1{,}4 (2{,}2) & 0{,}519 & $-$0{,}9 (2{,}4) & 0{,}720 \\
    \midrule
    Gemeinsamer Test & \multicolumn{2}{c}{F(2; 2726{,}6)\,=\,0{,}57, p\,=\,0{,}563}
     & \multicolumn{2}{c}{F(2; 2726{,}9)\,=\,2{,}48, p\,=\,0{,}084} \\
    \bottomrule
  \end{tabularx}

  \vspace{0.75ex}
  \parbox{\textwidth}{\footnotesize
    \textit{Anmerkungen:} Durchschnittliche marginale Effekte einer Erh\"ohung
    des Kreisangebots um eine Standardabweichung auf die Wahrscheinlichkeit
    bezahlter Arbeit, in Prozentpunkten, aus den Modellen der
    Tabelle~\ref{tab:ext-asyl-erwerb} (\texttt{mimrgns}, Poolung nach Rubins
    Regeln). Kontrolliert ist die Identit\"at
    AME(Gruppe)\,=\,$b$[Angebot]\,$+$\,$b$[Angebot\,$\times$\,Gruppe];
    gr\"o\ss{}te Abweichung zwischen beiden Rechenwegen 0{,}050~p.p.
    \\[0.4ex]
    \textbf{Inferenzregel, vorab festgelegt} (Designdokument, Abschnitt~3e):
    Prim\"ar ist der gemeinsame Test beider Interaktionsterme, einzelne
    Kontraste sind deskriptiv. \emph{Keine} der beiden Spezifikationen
    erreicht p\,$<$\,0{,}05. Der einzelne signifikante Kontrast in
    Modell~E1.2 darf daher nicht als Befund berichtet werden.
    \\[0.4ex]
    \textbf{Statistische Power, ebenfalls vorab festgehalten:} Bei
    Gruppenanteilen von 39/38/23~Prozent liegen die Standardfehler der
    Subgruppen etwa beim 1{,}6- bis 2{,}1-Fachen des gepoolten Werts von
    1{,}19. Ein Unterschied unter rund 4~Prozentpunkten ist auf dieser
    Ergebnisstufe nicht nachweisbar. Die Aussage lautet deshalb nicht
    \glqq es gibt keinen Effekt\grqq{}, sondern: ein Effekt dieser
    Gr\"o\ss{}enordnung h\"atte hier nicht entdeckt werden k\"onnen.
    Die Erweiterung wird auf der Mechanismenstufe entschieden
    (Tabelle~\ref{tab:ext-asyl-mechanismen}).}
\end{table}

\begin{table}[htbp]
  \centering
  \caption{Modell~E1.1 im Vergleich zu Tabelle~S6, Modell~1.1.9 der
           Publikation}
  \label{tab:ext-asyl-s6}
  \footnotesize
  \setlength{\tabcolsep}{4pt}%
  \begin{tabularx}{\textwidth}{@{}Zrrr@{}}
    \toprule
     & Publikation & Replikation & Abweichung \\
    \midrule
    Kreisangebot, std. (d.\,h.\ f\"ur \glqq keine Entscheidung\grqq{}) & 3{,}98$^{*}$ (1{,}55) & 3{,}97$^{*}$ (1{,}55) & 0{,}01 \\
    \quad $\times$ Anerkannt & $-$1{,}30 (1{,}90) & $-$1{,}28 (1{,}90) & 0{,}02 \\
    \quad $\times$ Abgelehnt & $-$2{,}58 (2{,}48) & $-$2{,}56 (2{,}47) & 0{,}02 \\
    \midrule
    N Beobachtungen & 5.473 & 5.467 & 6 \\
    N Personen & 2.732 & 2.730 & 2 \\
    \bottomrule
  \end{tabularx}

  \vspace{0.75ex}
  \parbox{\textwidth}{\footnotesize
    \textit{Anmerkungen:} Publizierte Werte aus dem Zusatzmaterial zu
    \textcite{Kanas:2023}, S.~12, Tabelle~S6, Modell~1.1.9; eigene Werte aus
    SOEPremote-Job~244835 (Modell~E1.1). Die Autorinnen sch\"atzen dort exakt
    diese Spezifikation, weshalb Modell~E1.1 zugleich eine Replikation von
    S6 ist. Alle drei Koeffizienten stimmen bis auf 0{,}02~p.p.
    \"uberein, die Standardfehler bis auf 0{,}01, und die Signifikanzniveaus
    vollst\"andig -- dieselbe Toleranz wie in den f\"unf Haupttabellen; die
    Differenz ist Monte-Carlo-Variation der multiplen Imputation.
    \\[0.4ex]
    \textbf{Vierter Fehler in der Publikation:} Tabelle~S6 weist
    5.473 Beobachtungen und 2.732 Personen aus, w\"ahrend die Tabellen~3
    bis~5 auf derselben Stichprobe und mit denselben Kontrollvariablen
    5.467 und 2.730 ausweisen. Die Nachsch\"atzung ergibt 5.467/2.730,
    passend zum Stichprobentrichter
    (Tabelle~\ref{tab:stichprobenselektion}). Die sechs zus\"atzlichen
    Personenjahre und zwei zus\"atzlichen Personen in S6 existieren nicht.
    Die Sch\"atzwerte sind davon unber\"uhrt.}
\end{table}

% Author  : Emre Suemer, 2026-09-01
%==============================================================================%

\makeatletter
\@ifundefined{NC@rewrite@Z}{%
  \newcolumntype{Z}{>{\raggedright\arraybackslash\hangindent=1.9em\hangafter=1}X}%
}{}
\makeatother

\begin{table}[htbp]
  \centering
  \caption{Kursangebot und Asylstatus: Wahrscheinlichkeit bezahlter Arbeit}
  \label{tab:ext-asyl-erwerb}
  \footnotesize
  \setlength{\tabcolsep}{4pt}%
  \begin{tabularx}{\textwidth}{@{}Zrr@{}}
    \toprule
     & Modell E1.1 & Modell E1.2 \\
     & Status im & Status bei \\
     & Erhebungsjahr & Erstbeobachtung \\
     & p.p. (SE) & p.p. (SE) \\
    \midrule
    \textbf{Kreisangebot an Integrationskursen, std.} & 3{,}97$^{*}$ (1{,}55) & 4{,}25$^{**}$ (1{,}31) \\
    \quad $\times$ Anerkannt & $-$1{,}28 (1{,}90) & $-$2{,}87 (2{,}03) \\
    \quad $\times$ Abgelehnt & $-$2{,}56 (2{,}47) & $-$5{,}13$^{*}$ (2{,}55) \\
    Aufenthaltsdauer, in Monaten (/12) & 10{,}73$^{**}$ (1{,}60) & 11{,}37$^{**}$ (1{,}60) \\
    Bildungsjahre vor der Migration & 0{,}56$^{**}$ (0{,}19) & 0{,}57$^{**}$ (0{,}20) \\
    Erwerbserfahrung vor der Migration & 4{,}64$^{+}$ (2{,}38) & 4{,}56$^{+}$ (2{,}39) \\
    Weiblich & $-$15{,}76$^{**}$ (2{,}16) & $-$15{,}69$^{**}$ (2{,}18) \\
    Kinder unter 16 Jahren im Haushalt & $-$11{,}02$^{**}$ (2{,}15) & $-$11{,}27$^{**}$ (2{,}18) \\
    Einreise mit der Familie & 5{,}33$^{**}$ (1{,}97) & 5{,}38$^{**}$ (2{,}01) \\
    Unterst\"utzung durch Familie/Verwandte in DE & 4{,}01 (2{,}83) & 4{,}10 (2{,}90) \\
    Alter bei der Erstbefragung & $-$0{,}35$^{**}$ (0{,}09) & $-$0{,}34$^{**}$ (0{,}09) \\
    Traumatische Erfahrungen & $-$0{,}24 (2{,}31) & $-$0{,}11 (2{,}33) \\
    Status des Asylantrags (Ref.: keine Entscheidung) & & \\
    \quad -- Anerkannt & 5{,}27$^{*}$ (2{,}61) & 5{,}21$^{*}$ (2{,}63) \\
    \quad -- Abgelehnt & 8{,}20$^{**}$ (2{,}76) & 4{,}39 (2{,}95) \\
    Kreisanteil Ausl\"ander:innen, std. & $-$1{,}12 (2{,}79) & $-$0{,}80 (2{,}83) \\
    Kreisbev\"olkerungsdichte, std. & $-$0{,}95 (3{,}61) & $-$1{,}05 (3{,}62) \\
    Kreisarbeitslosenquote, std. & $-$4{,}43$^{+}$ (2{,}43) & $-$3{,}99 (2{,}48) \\
    Kreissorgen \"uber Zuwanderung, std. & 0{,}33 (1{,}19) & 0{,}30 (1{,}21) \\
    Kreisanteil CDU/CSU-W\"ahler:innen, std. & $-$2{,}01 (2{,}20) & $-$1{,}65 (2{,}21) \\
    Konstante & $-$2{,}82 (6{,}31) & $-$4{,}16 (6{,}36) \\
    \midrule
    \multicolumn{3}{@{}l}{\textit{Fixe Effekte}} \\
    \quad Erhebungsjahr, Erhebungsstichprobe & Ja & Ja \\
    \quad Bundesland, Herkunftsland (aggregiert) & Ja & Ja \\
    \midrule
    Gemeinsamer Test beider Interaktionen & F(2; 2726{,}6)\,=\,0{,}57 & F(2; 2726{,}9)\,=\,2{,}48 \\
    \quad p & 0{,}563 & 0{,}084 \\
    \midrule
    N Beobachtungen & 5.467 & 5.467 \\
    N Personen & 2.730 & 2.730 \\
    N Imputationen & 20 & 20 \\
    R\textsuperscript{2} (adj.) & 0{,}2137 & 0{,}2121 \\
    \bottomrule
  \end{tabularx}

  \vspace{0.75ex}
  \parbox{\textwidth}{\footnotesize
    \textit{Anmerkungen:} Eigene Berechnung, IAB-BAMF-SOEP-Befragung von
    Gefl\"uchteten 2016--2019, SOEP-Core v36 (Remote Edition), gewichtet.  $^{+}$~p\,$<$\,0{,}1, $^{*}$~p\,$<$\,0{,}05, $^{**}$~p\,$<$\,0{,}01
    (zweiseitig).}
    
	%   Spezifikation wie Modell~1.1 in Tabelle~\ref{tab:erwerbstaetigkeit},
	%   erg\"anzt um die Interaktion des Kursangebots mit dem Asylstatus.
	%   Lineare Wahrscheinlichkeitsmodelle, durchschnittliche marginale Effekte in
	%   Prozentpunkten; robuste, auf Personenebene geclusterte Standardfehler in
	%   Klammern. Referenzkategorie des Asylstatus ist \glqq keine
	%   Entscheidung\grqq{}; die Kategorien sind gegen\"uber der Publikation
	%   korrigiert (vgl.\ Tabelle~\ref{tab:erwerbstaetigkeit}).
	%   Modell~E1.1 misst den Status im jeweiligen Erhebungsjahr, Modell~E1.2 bei
	%   der Erstbeobachtung der Person -- Letzteres ist konzeptionell sauberer,
	%   weil Asylentscheidungen erst \emph{nach} der Kreiszuweisung ergehen
	%   (Tabelle~\ref{tab:ext-asyl-uebergang}). In Spalte~E1.2 beziehen sich
	%   sowohl die Interaktions- als auch die Haupteffektzeilen des Asylstatus
	%   entsprechend auf den Status bei Erstbeobachtung.
    
\end{table}

\begin{table}[htbp]
  \centering
  \caption{Marginaler Effekt des Kursangebots auf die Erwerbst\"atigkeit,
           nach Asylstatus}
  \label{tab:ext-asyl-erwerb-ame}
  \footnotesize
  \setlength{\tabcolsep}{4pt}%
  \begin{tabularx}{\textwidth}{@{}Zrrrr@{}}
    \toprule
     & \multicolumn{2}{c}{Modell E1.1: Erhebungsjahr}
     & \multicolumn{2}{c}{Modell E1.2: Erstbeobachtung} \\
    \cmidrule(lr){2-3}\cmidrule(lr){4-5}
    Status des Asylantrags & p.p. (SE) & p & p.p. (SE) & p \\
    \midrule
    Keine Entscheidung & 4{,}0$^{*}$ (1{,}5) & 0{,}010 & 4{,}3$^{**}$ (1{,}3) & 0{,}001 \\
    Anerkannt & 2{,}7$^{+}$ (1{,}5) & 0{,}078 & 1{,}4 (1{,}9) & 0{,}467 \\
    Abgelehnt & 1{,}4 (2{,}2) & 0{,}519 & $-$0{,}9 (2{,}4) & 0{,}720 \\
    \midrule
    Gemeinsamer Test & \multicolumn{2}{c}{F(2; 2726{,}6)\,=\,0{,}57, p\,=\,0{,}563}
     & \multicolumn{2}{c}{F(2; 2726{,}9)\,=\,2{,}48, p\,=\,0{,}084} \\
    \bottomrule
  \end{tabularx}

  \vspace{0.75ex}
  \parbox{\textwidth}{\footnotesize
    \textit{Anmerkungen:} Durchschnittliche marginale Effekte einer Erh\"ohung
    des Kreisangebots um eine Standardabweichung auf die Wahrscheinlichkeit
    bezahlter Arbeit, in Prozentpunkten, aus den Modellen der
    Tabelle~\ref{tab:ext-asyl-erwerb} (\texttt{mimrgns}, Poolung nach Rubins
    Regeln). Kontrolliert ist die Identit\"at
    AME(Gruppe)\,=\,$b$[Angebot]\,$+$\,$b$[Angebot\,$\times$\,Gruppe];
    gr\"o\ss{}te Abweichung zwischen beiden Rechenwegen 0{,}050~p.p.
    \\[0.4ex]
    \textbf{Inferenzregel, vorab festgelegt} (Designdokument, Abschnitt~3e):
    Prim\"ar ist der gemeinsame Test beider Interaktionsterme, einzelne
    Kontraste sind deskriptiv. \emph{Keine} der beiden Spezifikationen
    erreicht p\,$<$\,0{,}05. Der einzelne signifikante Kontrast in
    Modell~E1.2 darf daher nicht als Befund berichtet werden.
    \\[0.4ex]
    \textbf{Statistische Power, ebenfalls vorab festgehalten:} Bei
    Gruppenanteilen von 39/38/23~Prozent liegen die Standardfehler der
    Subgruppen etwa beim 1{,}6- bis 2{,}1-Fachen des gepoolten Werts von
    1{,}19. Ein Unterschied unter rund 4~Prozentpunkten ist auf dieser
    Ergebnisstufe nicht nachweisbar. Die Aussage lautet deshalb nicht
    \glqq es gibt keinen Effekt\grqq{}, sondern: ein Effekt dieser
    Gr\"o\ss{}enordnung h\"atte hier nicht entdeckt werden k\"onnen.
    Die Erweiterung wird auf der Mechanismenstufe entschieden
    (Tabelle~\ref{tab:ext-asyl-mechanismen}).}
\end{table}

\begin{table}[htbp]
  \centering
  \caption{Modell~E1.1 im Vergleich zu Tabelle~S6, Modell~1.1.9 der
           Publikation}
  \label{tab:ext-asyl-s6}
  \footnotesize
  \setlength{\tabcolsep}{4pt}%
  \begin{tabularx}{\textwidth}{@{}Zrrr@{}}
    \toprule
     & Publikation & Replikation & Abweichung \\
    \midrule
    Kreisangebot, std. (d.\,h.\ f\"ur \glqq keine Entscheidung\grqq{}) & 3{,}98$^{*}$ (1{,}55) & 3{,}97$^{*}$ (1{,}55) & 0{,}01 \\
    \quad $\times$ Anerkannt & $-$1{,}30 (1{,}90) & $-$1{,}28 (1{,}90) & 0{,}02 \\
    \quad $\times$ Abgelehnt & $-$2{,}58 (2{,}48) & $-$2{,}56 (2{,}47) & 0{,}02 \\
    \midrule
    N Beobachtungen & 5.473 & 5.467 & 6 \\
    N Personen & 2.732 & 2.730 & 2 \\
    \bottomrule
  \end{tabularx}

  \vspace{0.75ex}
  \parbox{\textwidth}{\footnotesize
    \textit{Anmerkungen:} Publizierte Werte aus dem Zusatzmaterial zu
    \textcite{Kanas:2023}, S.~12, Tabelle~S6, Modell~1.1.9; eigene Werte aus
    SOEPremote-Job~244835 (Modell~E1.1). Die Autorinnen sch\"atzen dort exakt
    diese Spezifikation, weshalb Modell~E1.1 zugleich eine Replikation von
    S6 ist. Alle drei Koeffizienten stimmen bis auf 0{,}02~p.p.
    \"uberein, die Standardfehler bis auf 0{,}01, und die Signifikanzniveaus
    vollst\"andig -- dieselbe Toleranz wie in den f\"unf Haupttabellen; die
    Differenz ist Monte-Carlo-Variation der multiplen Imputation.
    \\[0.4ex]
    \textbf{Vierter Fehler in der Publikation:} Tabelle~S6 weist
    5.473 Beobachtungen und 2.732 Personen aus, w\"ahrend die Tabellen~3
    bis~5 auf derselben Stichprobe und mit denselben Kontrollvariablen
    5.467 und 2.730 ausweisen. Die Nachsch\"atzung ergibt 5.467/2.730,
    passend zum Stichprobentrichter
    (Tabelle~\ref{tab:stichprobenselektion}). Die sechs zus\"atzlichen
    Personenjahre und zwei zus\"atzlichen Personen in S6 existieren nicht.
    Die Sch\"atzwerte sind davon unber\"uhrt.}
\end{table}
